\chapter{Происхождение коэффициента кубического связного оператора}
\label{ch:v8-baryon-cubic-trace-parent-action-coefficient-origin-no-go-gate}

Предыдущая глава построила канонический связный оператор $W_3$, но оставила
свободным коэффициент $\lambda_3$. Проверим, выбирают ли его уже принятые
квадратичная форма Дирихле, ограниченность действия и стационарность.

\section{Точный луч с ненулевым кубическим следом}

В центрированном полном кадре выберем
\begin{equation}
 A=\widehat F_0+\widehat F_{40}.
 \label{eq:v8-baryon-cubic-coefficient-test-ray}
\end{equation}
Его первые три ненулевых следовых момента равны
\begin{equation}
 \Tr A^2=38,\qquad \Tr A^3=-3,\qquad \Tr A^4=134.
 \label{eq:v8-baryon-cubic-coefficient-ray-moments}
\end{equation}
Значит, кубический носитель не является формально нулевым даже на
одномерном рациональном срезе.

\section{Квадратичный родитель и ограниченность}

Самое общее следовое завершение до четвёртой степени на этом луче имеет вид
\begin{equation}
 S(t)=38\alpha t^2-3\lambda_3t^3+134\beta t^4,
 \qquad \alpha>0,\quad\beta>0.
 \label{eq:v8-baryon-cubic-coefficient-ray-action}
\end{equation}
Исходная квадратичная форма соответствует $\lambda_3=\beta=0$ и имеет
нулевую третью производную. Добавленный кубический член удовлетворяет
\begin{equation}
 S'''(0)=-18\lambda_3.
 \label{eq:v8-baryon-cubic-coefficient-third-derivative}
\end{equation}
Поэтому он не возникает из гессиана или квантово-марковской формы второго
порядка.

Чистый ненулевой кубический полином не ограничен снизу: смена знака $t$
меняет знак старшего члена. При $\beta>0$ квартика ограничивает действие
снизу при любом вещественном $\lambda_3$. Следовательно, ограниченность
отсекает чистый кубический родитель, но не выбирает коэффициент в
стабилизированном действии.

\section{Стационарность также оставляет семейство}

Если условно потребовать стационарность точки $t=1$, то
\begin{equation}
 S'(1)=0
 \quad\Longleftrightarrow\quad
 \lambda_3=\frac{76\alpha+536\beta}{9}.
 \label{eq:v8-baryon-cubic-coefficient-stationarity}
\end{equation}
После этой подстановки радиальная вторая производная равна
\begin{equation}
 S''(1)=-76\alpha+536\beta,
 \qquad
 S''(1)>0\Longleftrightarrow\beta>\frac{19}{134}\alpha.
 \label{eq:v8-baryon-cubic-coefficient-stability}
\end{equation}
Это открытый конус пар $(\alpha,\beta)$, а не единственная точка. Даже
безразмерная комбинация
\begin{equation}
 \frac{\lambda_3^2}{\alpha\beta}
 \label{eq:v8-baryon-cubic-coefficient-shape-ratio}
\end{equation}
не определяется следовой метрикой.

\begin{theorem}[Запрет вывода коэффициента из квадратичного родителя]
Текущая квадратичная форма Дирихле имеет нулевую третью вариацию и не
порождает $\lambda_3W_3$. Ограниченность снизу требует стабилизирующего
чётного члена, но допускает любое $\lambda_3$. Стационарность и локальная
устойчивость связывают коэффициент с двумя новыми параметрами и также не
выбирают его единственно.
\end{theorem}

\section{Единственный допустимый маршрут продолжения}

Опора $TTG$ совпадает по типу с перекрёстным членом между кривизной
калибровочной связности и квадратом поля переноса в стандартной
суперсвязностной кривизне. Но конечный внутренний след не содержит
пространственно-временную степень формы, оператор Ходжа, производную и общий
коэффициент кинетического действия. Поэтому это совпадение является
типизированной зацепкой, а не выводом $\lambda_3$.

\begin{nogocriterion}[Условие повторного открытия]
Новый свободный член $\lambda_3W_3$ добавлять запрещено. Следующий гейт
должен построить полную пространственно-временную супер-кривизну, вычислить
её единый нормированный квадрат и спроецировать его кубическую часть на
$W_3$. Только ненулевой коэффициент, полученный вместе с квадратичным и
квартичным членами из одного следа, повышает оператор до родительского.
\end{nogocriterion}

\begin{protocol}
\begin{enumerate}
 \item поле вычислений:
       \textbf{$\mathbb Q(\alpha,\beta,\lambda_3,t)$};
 \item моменты луча: \textbf{$(38,-3,134)$};
 \item третья вариация квадратичного родителя: \textbf{$0$};
 \item чистый кубический родитель ограничен снизу: \textbf{нет};
 \item положительная квартика допускает любой $\lambda_3$: \textbf{да};
 \item стационарность выбирает число: \textbf{нет};
 \item область устойчивости: \textbf{$\beta>19\alpha/134$};
 \item кубический носитель запрещён: \textbf{нет};
 \item коэффициент текущим действием выведен: \textbf{нет};
 \item следующий гейт:
       \textbf{проекция полной супер-кривизны на $W_3$}.
\end{enumerate}
\end{protocol}

Машинный сертификат сохранён в
\path{s2t_v8_baryon_cubic_trace_parent_action_coefficient_origin_no_go_gate_results.json}.